% !TEX root = ../main.tex
\begin{ouverture}[son]{21}{Intensité sonore et atténuation}{Partie 4 · Ondes et signaux}
{Intensité sonore\enspace·\enspace Niveau d'intensité sonore\enspace·\enspace Décibel\enspace·\enspace Atténuation géométrique et par absorption}
\objectif{Exploiter l'expression du niveau d'intensité sonore d'un signal.}
\objectif{Utiliser la fonction logarithme décimal et sa fonction réciproque.}
\objectif{Calculer le niveau sonore produit par plusieurs sources et une atténuation en décibels.}
\objectif{Illustrer et distinguer l'atténuation géométrique et l'atténuation par absorption.}
\objectif{Calculer l'évolution d'un niveau sonore avec la distance à l'aide d'un programme Python.}
\prerequis{Signal sonore, fréquence, célérité (classe de seconde).}
\prerequis{Puissance par unité de surface d'une onde sonore (classe de première).}
\prerequis{Fonction logarithme décimal (mathématiques).}
\end{ouverture}
\essentiel
\essSection{1}{L'intensité sonore}
\begin{definition}[Intensité sonore]
L'\textbf{intensité sonore} $I$ est la puissance transportée par l'onde sonore par unité de surface perpendiculaire
à la direction de propagation : $I=\dfrac{P}{S}$, en \unit{W.m^{-2}}. L'oreille humaine perçoit, à \qty{1}{kHz}, des
intensités comprises entre le \textbf{seuil d'audibilité} $I_0=\qty{1,0e-12}{W.m^{-2}}$ et le seuil de douleur, voisin
de \qty{1}{W.m^{-2}} : douze ordres de grandeur.
\end{definition}
\begin{propriete}[Source ponctuelle]
Une source ponctuelle qui émet la puissance sonore $P$ de façon identique dans toutes les directions répartit cette
puissance sur des sphères de surface $4\pi d^2$ : à la distance $d$,
\[I=\frac{P}{4\pi d^2}.\]
L'intensité est divisée par 4 lorsque la distance double.
\end{propriete}
\schema{21-sphere}
\essSection{2}{Le niveau d'intensité sonore}
\begin{definition}[Niveau d'intensité sonore]
Le \textbf{niveau d'intensité sonore} $L$, exprimé en \textbf{décibels} (dB), est défini par
\[L=10\log\left(\frac{I}{I_0}\right)\qquad\text{et, réciproquement,}\qquad I=I_0\times10^{L/10},\]
où $\log$ désigne le logarithme décimal. Le seuil d'audibilité correspond à \qty{0}{dB}, le seuil de douleur à
\qty{120}{dB}.
\end{definition}
\begin{propriete}[Additionner des sources sonores]
Les intensités sonores de sources indépendantes s'ajoutent ; les niveaux, eux, ne s'ajoutent pas. Multiplier
l'intensité par 10 augmente le niveau de \qty{10}{dB} ; la multiplier par 2 l'augmente de $10\log2\approx\qty{3}{dB}$ :
deux sources identiques produisent \qty{3}{dB} de plus qu'une seule.
\end{propriete}
\schema{21-echelle}
\begin{savoirfaire}[Calculer avec des décibels]
Convertir chaque niveau en intensité, additionner les intensités, puis revenir au niveau. Trois machines produisant
chacune \qty{80}{dB} en un point y apportent chacune $I=10^{-12}\times10^{8}=\qty{1,0e-4}{W.m^{-2}}$ ; ensemble,
$I=\qty{3,0e-4}{W.m^{-2}}$, d'où $L=10\log(\num{3.0e8})=\qty{84,8}{dB}$, et non \qty{240}{dB} !
\end{savoirfaire}
\essSection{3}{L'atténuation}
\begin{definition}[Atténuation]
L'\textbf{atténuation} entre deux points, ou due à un dispositif, est la diminution du niveau d'intensité sonore :
$A=L_1-L_2$, en décibels.
\end{definition}
\begin{propriete}[Atténuation géométrique et atténuation par absorption]
\begin{itemize}
\item L'\textbf{atténuation géométrique} résulte de la répartition de la puissance sur une surface de plus en plus grande.
Pour une source ponctuelle, $L_2=L_1-20\log\dfrac{d_2}{d_1}$ : le niveau baisse de \qty{6}{dB} quand la distance double.
\item L'\textbf{atténuation par absorption} résulte de la conversion d'une partie de l'énergie de l'onde en énergie
interne du milieu traversé (air, mur, mousse, protections auditives). Elle dépend du matériau, de son épaisseur et de
la fréquence ; dans l'air, les sons aigus sont plus absorbés que les sons graves.
\end{itemize}
\end{propriete}
\schema{21-attenuation}
\begin{codecote}[0.58]
\begin{pythoncol}
from math import log10, pi
P, I0 = 0.10, 1e-12            # W, W/m²
for d in [1, 2, 4, 8, 16]:     # m
    I = P / (4 * pi * d**2)
    L = 10 * log10(I / I0)
    print(d, round(L, 1))
\end{pythoncol}
\begin{consolecol}
1 99.0
2 93.0
4 87.0
8 80.9
16 74.9
\end{consolecol}
\cote
\fontsize{9.5}{12.5}\selectfont
\capacite{Capacité numérique}\par\smallskip
Le programme calcule l'intensité, puis le niveau sonore produits par une source ponctuelle de \qty{0,10}{W} à
différentes distances. Chaque doublement de la distance divise l'intensité par 4 et abaisse le niveau de
\qty{6,0}{dB}, aux arrondis près.
\end{codecote}
\begin{attention}
On n'additionne jamais des décibels : on additionne les intensités. Par ailleurs, la sensation auditive dépend aussi de
la fréquence ; les sonomètres affichent souvent un niveau pondéré, en dB(A), qui tient compte de la sensibilité de
l'oreille.
\end{attention}
\begin{experience}[Mesurer une atténuation]
Face à un haut-parleur alimenté par un générateur basse fréquence, on relève le niveau sonore avec un sonomètre à
différentes distances : en plein air, il diminue d'environ \qty{6}{dB} à chaque doublement de la distance. On interpose
ensuite des plaques de matériaux différents (bois, mousse, polystyrène) entre la source et le sonomètre placé à
distance fixe : la différence de niveau mesure l'atténuation par absorption, qui augmente avec l'épaisseur et la
fréquence.
\end{experience}
\begin{rediger}
Écrire la relation utilisée, $L=10\log(I/I_0)$ ou $I=I_0\times10^{L/10}$, et préciser l'unité de chaque grandeur. Pour
plusieurs sources, rédiger en trois étapes : « intensité de chaque source », « intensité totale », « niveau total ».
Préciser la nature de l'atténuation : géométrique ou par absorption.
\end{rediger}
\exercices

\rubrique{Découverte}
\begin{blocLivre}{ex-21-1}
\begin{exercice}[1]{Du niveau à l'intensité}
\begin{donnees}
$I_0=\qty{1,0e-12}{W.m^{-2}}$ ; $L=10\log(I/I_0)$.
\end{donnees}
\begin{questions}
\item Calculer l'intensité sonore d'une conversation de \qty{60}{dB}.
\item Calculer le niveau sonore d'une rue animée où $I=\qty{3,2e-5}{W.m^{-2}}$.
\item À partir de quel niveau sonore une exposition prolongée devient-elle dangereuse ? Calculer l'intensité
correspondante, sachant que ce seuil vaut \qty{85}{dB}.
\item Combien de fois l'intensité d'un concert à \qty{100}{dB} est-elle plus grande que celle d'une rue à \qty{70}{dB} ?
\end{questions}
\end{exercice}
\end{blocLivre}
\begin{blocLivre}{ex-21-2}
\begin{exercice}[1]{Un, deux, dix violons}
Un violoniste produit, à l'endroit où se trouve un auditeur, un niveau sonore de \qty{70}{dB}. Tous les violonistes
d'un orchestre sont supposés produire le même niveau en ce point.
\begin{questions}
\item Calculer le niveau sonore produit par deux violons. Le niveau a-t-il doublé ?
\item Calculer le niveau sonore produit par dix violons.
\item Combien de violons faut-il pour atteindre \qty{85}{dB} ?
\item Pourquoi utilise-t-on une échelle logarithmique pour décrire les sons ?
\end{questions}
\end{exercice}
\end{blocLivre}
\begin{blocLivre}{ex-21-3}
\begin{exercice}[1]{S'éloigner d'une enceinte}
À \qty{2,0}{m} d'une enceinte, assimilée à une source ponctuelle, le niveau sonore vaut \qty{100}{dB}. On néglige
l'absorption par l'air.
\begin{donnees}
$I_0=\qty{1,0e-12}{W.m^{-2}}$.
\end{donnees}
\begin{questions}
\item Calculer le niveau sonore à \qty{4,0}{m}, puis à \qty{20}{m}.
\item Quel type d'atténuation est en jeu ?
\item Calculer la puissance sonore émise par l'enceinte.
\item À quelle distance le niveau n'est-il plus que de \qty{70}{dB} ?
\end{questions}
\end{exercice}
\end{blocLivre}
\begin{blocLivre}{ex-21-4}
\begin{exercice}[1]{Les protections auditives}
Sur un chantier, un ouvrier est exposé à \qty{105}{dB}. Il porte des bouchons d'oreille dont l'atténuation vaut
\qty{25}{dB}. On admet qu'à \qty{85}{dB} l'oreille supporte 8~heures d'exposition par jour, durée divisée par deux pour
chaque hausse de \qty{3}{dB}.
\begin{questions}
\item Quel type d'atténuation les bouchons réalisent-ils ? Quel niveau sonore l'oreille reçoit-elle ?
\item Quelle fraction de l'intensité sonore les bouchons laissent-ils passer ?
\item Quelle est la durée d'exposition admissible à \qty{94}{dB} ?
\item Combien de temps l'ouvrier pourrait-il travailler sans bouchons avant d'atteindre sa dose quotidienne ?
\end{questions}
\end{exercice}
\end{blocLivre}
\begin{blocLivre}{ex-21-5}
\begin{exercice}[1]{Logarithme et puissances de dix}
\begin{questions}
\item Calculer sans calculatrice $\log(1000)$, $\log(\num{1e-12})$ et $\log(1)$.
\item À la calculatrice, calculer $\log2$ et $10^{0{,}5}$, puis en déduire l'effet sur le niveau sonore d'une
intensité doublée.
\item Montrer que $L=10\log\dfrac{I}{I_0}$ et $I=I_0\times10^{L/10}$ sont deux écritures de la même relation.
\item Les intensités audibles s'étendent de \qty{1e-12}{W.m^{-2}} à \qty{1}{W.m^{-2}}. Combien d'ordres de grandeur cela
représente-t-il ? À quel intervalle de niveaux cela correspond-il ?
\end{questions}
\end{exercice}
\end{blocLivre}
\rubrique{Entraînement}
\begin{blocLivre}{ex-21-6}
\begin{exercice}[2]{Au pied de la scène}
Lors d'un concert en plein air, le niveau sonore vaut \qty{110}{dB} à \qty{5,0}{m} des enceintes, assimilées à une source
ponctuelle.
\begin{questions}
\item Calculer le niveau sonore pour un spectateur placé à \qty{20}{m}, puis à \qty{40}{m}.
\item La réglementation impose de ne pas dépasser \qty{102}{dB} en moyenne là où se trouve le public. À quelle distance
minimale des enceintes faut-il placer la barrière ?
\item Un spectateur placé à \qty{20}{m} porte des bouchons d'atténuation \qty{20}{dB}. Quel niveau perçoit-il ?
\item Pourquoi le son entendu de très loin paraît-il plus grave ?
\end{questions}
\end{exercice}
\end{blocLivre}
\begin{blocLivre}{ex-21-7}
\begin{exercice}[2]{Le brouhaha de la cantine}
Dans une cantine, chaque conversation produit un niveau sonore de \qty{55}{dB} au centre de la salle. On néglige les
autres sources.
\begin{questions}
\item Calculer le niveau sonore au centre de la salle lorsque quinze conversations ont lieu en même temps.
\item Même question pour trente conversations. Commenter.
\item De combien le niveau augmente-t-il lorsqu'une seizième conversation s'ajoute aux quinze premières ?
\item Pour se faire entendre, chacun parle plus fort : expliquer pourquoi le niveau sonore d'une cantine a tendance à
augmenter de lui-même.
\end{questions}
\end{exercice}
\end{blocLivre}
\begin{blocLivre}{ex-21-8}
\begin{exercice}[2]{Programmer l'atténuation}
On reprend le programme de l'essentiel.
\begin{questions}
\item \capacite{Capacité numérique}Que représentent \code{P} et \code{I0} ? Quelle relation la ligne \code{I = P / (4 * pi * d**2)}
traduit-elle ?
\item Vérifier, à partir des valeurs affichées, que le niveau diminue de \qty{6}{dB} à chaque doublement de la distance.
\item L'air absorbe un son de \qty{1}{kHz} à raison d'environ \qty{5}{dB} par kilomètre. Modifier la ligne de calcul de
\code{L} pour en tenir compte. Ce terme est-il important à \qty{16}{m} ?
\item Calculer le niveau à \qty{1000}{m}, sans puis avec absorption.
\end{questions}
\end{exercice}
\end{blocLivre}
\begin{blocLivre}{ex-21-9}
\begin{exercice}[2]{Choisir ses fenêtres}
Une chambre donne sur une rue où le niveau sonore vaut \qty{72}{dB} devant la façade. Une fenêtre en simple vitrage
apporte une atténuation de \qty{25}{dB} ; une fenêtre à double vitrage acoustique, \qty{35}{dB}.
\begin{questions}
\item Quel type d'atténuation les vitrages réalisent-ils ? Calculer le niveau sonore dans la chambre dans les deux cas.
\item Combien de fois l'intensité transmise par le double vitrage est-elle plus faible que celle transmise par le simple
vitrage ?
\item Pour bien dormir, on recommande un niveau inférieur à \qty{35}{dB}. Quelle atténuation minimale faut-il ?
\item La fenêtre entrouverte n'atténue plus que de \qty{10}{dB}. Quelle fraction de l'intensité laisse-t-elle passer ?
\end{questions}
\end{exercice}
\end{blocLivre}
\begin{blocLivre}{ex-21-10}
\begin{exercice}[2]{Le grondement lointain du tonnerre}
Un éclair produit, à \qty{1}{m}, un niveau sonore de \qty{100}{dB} à toutes les fréquences audibles (valeur arbitraire
choisie pour comparer). L'air absorbe les sons d'autant plus que leur fréquence est élevée ; valeurs approximatives de
l'absorption :
\begin{center}
\begin{tabular}{lcccc}
\toprule
Fréquence (Hz) & 500 & \num{1000} & \num{4000} & \num{8000} \\
Absorption (dB par km) & 2,8 & 4,7 & 24 & 83 \\
\bottomrule
\end{tabular}
\end{center}
\begin{questions}
\item Calculer l'atténuation géométrique entre \qty{1}{m} et \qty{1}{km}, l'éclair étant assimilé à une source ponctuelle.
\item Calculer, pour chaque fréquence, le niveau sonore à \qty{1}{km} en tenant compte des deux atténuations.
\item Quelles fréquences sont encore audibles à \qty{1}{km} ?
\item Expliquer pourquoi le tonnerre lointain gronde sourdement alors que le tonnerre proche claque.
\end{questions}
\end{exercice}
\end{blocLivre}
\begin{blocLivre}{ex-21-11}
\begin{exercice}[2]{Source ponctuelle ou source linéaire}
Une autoroute très fréquentée se comporte comme une source linéaire : la puissance sonore se répartit sur des cylindres
de surface $2\pi d\,\ell$, où $\ell$ est la longueur considérée. À \qty{25}{m} de l'autoroute, le niveau sonore vaut
\qty{75}{dB}.
\begin{questions}
\item Montrer que, pour une source linéaire, l'intensité est inversement proportionnelle à la distance $d$.
\item En déduire que le niveau baisse de \qty{3}{dB} quand la distance double.
\item Calculer le niveau à \qty{100}{m} de l'autoroute. Quelle valeur aurait-on obtenue pour une source ponctuelle ?
\item À quelle distance de l'autoroute le niveau descend-il à \qty{60}{dB} ? Pourquoi construit-on des murs antibruit ?
\end{questions}
\end{exercice}
\end{blocLivre}
\rubrique{Objectif bac}
\begin{blocLivre}{ex-21-12}
\begin{exercice}[3]{Protéger le public d'un concert}
\begin{docu}{1}{La réglementation}
Depuis 2017, dans les lieux diffusant de la musique amplifiée, le niveau sonore moyen ne doit dépasser
\qty{102}{dB} sur 15~minutes en aucun endroit accessible au public.
\end{docu}
\begin{docu}{2}{Le dispositif}
Deux enceintes identiques, placées de part et d'autre de la scène, produisent chacune \qty{110}{dB} à \qty{4,0}{m}. On les
assimile à des sources ponctuelles et l'on considère un spectateur situé dans l'axe de la scène, à égale distance des
deux enceintes.
\end{docu}
\begin{questions}
\item Avec une seule enceinte en fonctionnement, à quelle distance minimale doit se trouver le public ?
\item Quel niveau produisent les deux enceintes ensemble pour un spectateur placé à cette distance ?
\item \initiative Avec les deux enceintes, à quelle distance minimale de chacune doit se trouver le public ?
\item Pourquoi les ingénieurs du son préfèrent-ils répartir de nombreuses petites enceintes plutôt que d'augmenter la
puissance de deux grosses ?
\end{questions}
\end{exercice}
\end{blocLivre}
\begin{blocLivre}{ex-21-13}
\begin{exercice}[3]{L'éolienne et ses voisins}
\begin{docu}{1}{Une éolienne bruyante}
Le niveau de puissance acoustique d'une éolienne vaut \qty{105}{dB} : sa puissance sonore $P$ vérifie
$10\log\dfrac{P}{P_0}=105$, avec $P_0=\qty{1e-12}{W}$. On l'assimile à une source ponctuelle et l'on néglige l'absorption
par l'air et le sol.
\end{docu}
\begin{docu}{2}{La réglementation de l'émergence}
La nuit, le bruit ambiant près d'une maison vaut \qty{35}{dB}. L'émergence, différence entre le niveau total avec
l'éolienne et le niveau ambiant sans elle, ne doit pas dépasser \qty{3}{dB} la nuit.
\end{docu}
\begin{donnees}
$I_0=\qty{1,0e-12}{W.m^{-2}}$.
\end{donnees}
\begin{questions}
\item Calculer la puissance sonore de l'éolienne, puis le niveau qu'elle produit seule à \qty{500}{m}.
\item Calculer le niveau total la nuit pour une maison située à \qty{500}{m}, puis l'émergence. La réglementation est-elle
respectée ?
\item \initiative À quelle distance minimale de l'éolienne doit se trouver la maison ?
\item L'absorption par l'air et par le sol a été négligée. Le résultat obtenu est-il trop exigeant ou pas assez ?
\end{questions}
\end{exercice}
\end{blocLivre}
\begin{blocLivre}{ex-21-14}
\begin{exercice}[3]{Choisir la fréquence d'une échographie}
\begin{docu}{1}{Des ultrasons absorbés}
En échographie, une sonde émet des ultrasons qui se réfléchissent sur les organes. Dans les tissus mous, l'atténuation
par absorption vaut environ \qty{0,5}{dB} par centimètre parcouru et par mégahertz de fréquence.
\end{docu}
\begin{docu}{2}{Contraintes de l'appareil}
L'échographe ne peut exploiter que les échos atténués de moins de \qty{60}{dB} sur l'aller-retour. Plus la fréquence est
élevée, plus l'image est fine.
\end{docu}
\begin{questions}
\item Calculer l'atténuation subie par une onde de \qty{5}{MHz} qui se réfléchit sur un organe situé à \qty{8}{cm} de
profondeur, aller et retour.
\item Quelle fraction de l'intensité émise revient alors à la sonde, en négligeant la réflexion partielle ?
\item \initiative Calculer la profondeur maximale explorable à \qty{7,5}{MHz}, puis à \qty{3,5}{MHz}.
\item Quelle fréquence choisir pour examiner la thyroïde, juste sous la peau ? pour un foie, à plus de \qty{10}{cm} ?
\end{questions}
\end{exercice}
\end{blocLivre}
\begin{blocLivre}{ex-21-15}
\begin{exercice}[3]{La dose de bruit quotidienne}
\begin{docu}{1}{Un niveau moyen}
L'exposition d'un travailleur se caractérise par un niveau moyen sur 8~heures : c'est le niveau d'un bruit constant qui
transporterait la même énergie sonore, $L_\mathrm{moy}=10\log\left(\dfrac{1}{8}\sum t_i\times10^{L_i/10}\right)$, les durées $t_i$
étant en heures. Le port de protections est obligatoire au-delà de \qty{85}{dB}.
\end{docu}
\begin{docu}{2}{La journée d'un menuisier}
1~heure de scie à \qty{94}{dB}, 3~heures de ponceuse à \qty{88}{dB} ; le reste du temps, l'exposition est négligeable.
\end{docu}
\begin{questions}
\item Expliquer pourquoi le niveau moyen fait intervenir les intensités, et non les niveaux.
\item Calculer le niveau moyen sur 8~heures du menuisier. Doit-il porter des protections ?
\item \initiative Le menuisier porte des bouchons de \qty{15}{dB} lorsqu'il utilise ses machines. Calculer son nouveau
niveau moyen.
\item Retrouver qu'une heure à \qty{94}{dB} transporte autant d'énergie sonore que huit heures à \qty{85}{dB}.
\end{questions}
\end{exercice}
\end{blocLivre}
